% ── Section: Abstract ────────────────────────────────────────────────────────
\begin{abstract}
Modern clinical AI systems store millions of high-dimensional vector
embeddings of patient records and must answer similarity search queries
in real time---while strictly enforcing fine-grained access control
policies mandated by regulations such as HIPAA.
Today, no vector database index handles this requirement efficiently
across the full spectrum of \emph{query selectivity}: the fraction of the
database a given user is permitted to see.
Post-filtering (retrieve candidates, then discard unauthorized results)
collapses to near-zero recall when access is restricted to $\leq 1\%$
of records.
Brute-force pre-filtering achieves exact results but degrades to
single-digit queries per second at broad access levels.
Multi-index architectures consume 8--64$\times$ more memory for realistic
role hierarchies.

We present \textbf{RBAC-HNSW}: a single modification to the
Hierarchical Navigable Small World (HNSW) graph-search algorithm that
enforces Role-Based Access Control (RBAC) at the node level, before the
expensive distance computation, while \emph{routing through} denied nodes
to preserve graph connectivity.
A 64-bit bitmask co-located with each vector adds only 0.24\% memory
overhead and costs $\approx\!1$ CPU cycle per access check.
On a 50{,}000-vector synthetic clinical dataset (BioBERT-scale,
768-dimensional), RBAC-HNSW achieves $\geq 0.90$ recall@10 across all
tested selectivity levels (0.02\%--40\%) and adds negligible memory
compared with the best multi-index competitor.
This paper presents the algorithm, a realistic HIPAA-motivated bitmask
schema for hospital access hierarchies, and an empirical evaluation
targeting the data-management community.
\end{abstract}
